\documentclass{article}
\usepackage{graphicx}
\usepackage{booktabs}

\title{PRM Framework: Process Ancestry-Based Security for HPC Environments}
\author{}
\date{}

\begin{document}

\maketitle

\begin{abstract}
Modern High-Performance Computing (HPC) environments face security challenges balancing protection with performance. This paper presents the PRM Framework, a process ancestry-based security solution for Linux that provides fine-grained control over system calls with minimal overhead (0.5-2\% vs. SELinux's 5-10\%).
\end{abstract}

\section{Introduction}
Security in HPC environments requires minimal overhead while addressing containerization risks. Traditional solutions like SELinux and AppArmor introduce performance penalties or lack context awareness.

\subsection{Contributions}
\begin{itemize}
\item Novel process ancestry-based security with minimal performance impact
\item Comprehensive architecture for efficient system call filtering
\item Performance benchmarks showing significant advantages over traditional solutions
\item Flexible policy framework for process ancestry-based rules
\end{itemize}

\section{Background}
\subsection{Linux Security Mechanisms}
Current Linux security includes DAC, MAC (SELinux/AppArmor), capabilities, namespaces, and seccomp.

\subsection{HPC Security Challenges}
HPC environments face unique challenges: performance sensitivity, complex workflows requiring privileges, multi-tenancy, and containerization risks.

\section{Related Work}
Prior work includes Systrace, Seccomp-BPF, SPEAKER, Docker Security, gVisor, and HPC-specific solutions.

\section{Methodology}
\subsection{Research Questions}
We address how to provide robust security without performance overhead, using process ancestry for context.

\subsection{Requirements}
Functional requirements include system call filtering based on process ancestry. Non-functional requirements include minimal overhead (<3\%) and compatibility with existing systems.

\section{PRM Framework Design}
\subsection{Architecture Overview}
The PRM Framework consists of:
\begin{itemize}
\item System Call Interception
\item Cache Mechanism
\item Filters
\item Verifier (Process Ancestry Checker)
\item PRM Table Rules
\item PROG Handlers
\end{itemize}

\subsection{Process Ancestry Analysis}
The framework extracts command names from current process, parent, and grandparent to make security decisions.

\section{Evaluation}
\subsection{Performance Comparison}
PRM Framework shows 0.5-2\% system call overhead compared to SELinux (5-10\%) and AppArmor (3-7\%).

\subsection{Performance Benchmarks}
Benchmarks on LINPACK, STREAM, IOR, and IMB-MPI show minimal impact on HPC workloads.

\subsection{Security Effectiveness}
The framework successfully prevented container escape attempts and privilege escalation.

\section{Discussion}
\subsection{Limitations}
Current limitations include configuration complexity and lack of integration with SIEM systems.

\section{Conclusion}
The PRM Framework provides a novel approach to Linux security well-suited for HPC environments, with future work focusing on automated rule generation and machine learning integration.

\end{document}